\section{What This Does Not Prove}
\label{sec:limitations}

The result that sigmoid transformers can exactly realize belief propagation
is formal and verified. The applications that follow range from
well-established to speculative. This section states clearly what the
theory does not prove.

\subsection*{The Theorem Is About Sigmoid Transformers}

Production transformers use SwiGLU or GeLU, not sigmoid activation. The
exact BP equivalence holds for sigmoid transformers under the constructive
weight mapping. For production transformers, the architecture is compatible
with BP-like computation and in synthetic experiments gradient descent
finds approximately correct BP solutions, but exact equivalence is not
established for arbitrary trained models.

\subsection*{Attention Is Not Literally Logical Conjunction}

The AND/OR framing is mechanistically motivated but is an interpretation,
not a theorem about arbitrary trained attention heads. Across three studied
model families (GPT-2 small, GPT-2 medium, Qwen~2.5 0.5B), a substantial
fraction of heads fall into the boolean-AND, continuous-OR, and hub
categories. The remainder exhibit context-dependent behavior not yet
characterized by the framework. Reproducibility and generalization to other
model families remain to be established.

\subsection*{The Factor Graph Is Not Directly Observable}

The implicit factor graph is encoded in learned weights, not explicitly
declared. Current SAEs cover a small fraction of the residual stream's
representational capacity. Factor graph edges have not been systematically
recovered. The claim that SAE features are natural candidates for factor
graph nodes is an empirical hypothesis with supporting evidence, not a
proven theorem.

\subsection*{What Would Falsify the Core Interpretation}

The BP interpretation would be significantly weakened if: boolean-AND heads
do not exhibit rank-1 Q/K structure under SAE analysis; SAE features fail
to organize into stable dependency structure; FFN updates are not locally
interpretable as belief updates; deeper layers do not improve iterative
reasoning behavior; or gradient descent does not find BP-like solutions on
clean synthetic tasks.